\documentclass[12pt]{article}
\usepackage[utf8]{inputenc}
\usepackage[T1]{fontenc}
\usepackage[margin=1in]{geometry}
\usepackage{setspace}
\usepackage{xcolor}
\usepackage[colorlinks=true, urlcolor=blue!60!black, linkcolor=blue!60!black]{hyperref}
\setlength{\parindent}{2em}
\setlength{\parskip}{1.25em}
\renewcommand{\baselinestretch}{1.0}
\begin{document}
\begin{center}
{\LARGE \textit{Friends of the Earth, Inc.\ v.\ Laidlaw Environmental Services (TOC), Inc.}, 528 U.S.\ 167 (2000)\par}
\vspace{0.8em}
{\large Case Brief --- Four-Page Edition\par}
\end{center}
\vspace{1em}

\section{Memory Jogger}

Environmental groups sued a polluter under the Clean Water Act for civil penalties. The Court held 7--2 that members who curtailed their use of a river had standing without proof of environmental harm, that penalties deterring ongoing violations provide redress, and that voluntary cessation moots a case only if it is absolutely clear the violations cannot reasonably be expected to recur (\href{https://supreme.justia.com/cases/federal/us/528/167/}{Justia}).

\section{Detailed Case Facts}

In 1986 Laidlaw bought a hazardous-waste incinerator in Roebuck, South Carolina, whose wastewater plant discharged into the North Tyger River. Its 1987 Clean Water Act discharge permit limited pollutants, including mercury, and Laidlaw committed nearly 500 permit violations between 1987 and 1995. In 1992 Friends of the Earth (FOE) and CLEAN, later joined by the Sierra Club, sent the required 60-day notice of a citizen suit. One day before the window closed, at Laidlaw's request, the state agency sued Laidlaw in state court on a complaint Laidlaw drafted, and settled the next day for \$100,000.

\section{Procedural History}

The district court refused to bar the citizen suit, finding the state settlement hasty and inadequate. After trial it found ``no demonstrated proof of harm to the environment,'' denied an injunction because Laidlaw had reached substantial compliance, and imposed a \$405,800 penalty to deter future violations. FOE appealed the amount; Laidlaw cross-appealed on standing. Relying on \textit{Steel Co.\ v.\ Citizens for a Better Environment}, the Fourth Circuit held the case moot because penalties payable to the Treasury could not redress FOE's injury. Meanwhile Laidlaw closed and dismantled the incinerator but kept its permit. The Supreme Court reversed.

\section{Judicial Votes}

\textbf{7--2}, reversed and remanded (January 12, 2000).
\begin{itemize}
\item \textbf{Majority:} Justice Ginsburg, joined by Chief Justice Rehnquist and Justices Stevens, O'Connor, Kennedy, Souter, and Breyer.
\item \textbf{Concurrences:} Justice Stevens; Justice Kennedy.
\item \textbf{Dissent:} Justice Scalia, joined by Justice Thomas.
\end{itemize}

\section{Holding}

Members' reasonable concerns that led them to curtail use of the river were injury in fact, even without proof of harm to the river. Civil penalties payable to the Treasury redress violations ongoing when suit is filed by deterring future noncompliance. And a defendant's voluntary cessation does not moot a case unless it carries the ``formidable burden'' of showing that recurrence cannot reasonably be expected---a question distinct from standing, left for remand.

\section{Analysis of Opinions}

\textbf{Majority (Ginsburg): injury.} An association has standing if its members do (\textit{Hunt}). The relevant showing is injury to the plaintiff, not to the environment. Unlike the ``some day'' intentions in \textit{Lujan}, members' affidavits showed they lived nearby and had stopped fishing, camping, wading, boating, and picnicking because of the discharges.

\textbf{Majority (Ginsburg): redress.} Standing is required for each form of relief, but penalties for violations ongoing at filing serve Congress's deterrent purpose, unlike the wholly past violations in \textit{Steel Co.}

\textbf{Majority (Ginsburg): mootness is not standing in a time frame.} The plaintiff bears the burden on standing at filing, but a defendant claiming voluntary cessation bears a heavy burden; and mootness exceptions exist to prevent strategic evasion after years of litigation. Whether compliance or closure made recurrence impossible was unresolved.

\textbf{Stevens (concurring).} Post-filing compliance cannot moot penalties already incurred, which serve punitive and deterrent ends.

\textbf{Kennedy (concurring).} Private suits for public fines raise ``difficult and fundamental questions'' under Article~II, not presented here.

\textbf{Scalia, joined by Thomas (dissenting).} Crediting subjective affidavits despite a finding of no harm makes injury a ``mere pleading requirement,'' and Treasury penalties give no personal redress; the Court put public enforcement power in private hands.

\textbf{Later developments.} \textit{Already v.\ Nike} (2013) and \textit{FBI v.\ Fikre} (2024) reaffirmed the formidable burden; \textit{TransUnion v.\ Ramirez} (2021) stands in tension with \textit{Laidlaw}'s view of injury.

\section{Examples: Future Applications}

\begin{enumerate}
\item \textbf{Same-Side: Post-Filing Compliance.} A plant that achieves compliance six months into a citizen suit cannot moot it without making recurrence absolutely clear.
\item \textbf{Same-Side: Birdwatchers.} Members who stopped visiting a marsh because of illegal dioxin discharges have standing even if scientists find no measurable harm.
\item \textbf{Opposite-Side: Violations Ended Before Suit.} If violations ceased before filing, nothing remains to deter; \textit{Steel Co.} and \textit{Gwaltney} defeat the suit.
\item \textbf{Opposite-Side: Undisclosed Credit Error.} A never-disclosed file error is a bare statutory violation without concrete harm under \textit{TransUnion}.
\item \textbf{Fence-Sitter: Closed Plant, Retained Permit.} A shuttered facility that keeps its permit, Laidlaw's own posture, raises a factual question that \textit{Fikre} suggests a bare assurance cannot settle.
\end{enumerate}

\section{Critique}

\textbf{Subjective fear.} Critics echo Scalia that separating injury to the plaintiff from injury to the environment risks manufacturing standing from anxiety over technical violations causing no ecological harm. \textbf{Article~II.} Kennedy's reservation remains open: private suits for public fines may trench on executive enforcement discretion. \textbf{Asymmetry.} \textit{Laidlaw} refutes the shorthand that mootness is standing in a time frame, shifting the burden to the defendant.

\section{Key Quotations}

\begin{quote}
``[A] defendant claiming that its voluntary compliance moots a case bears the formidable burden of showing that it is absolutely clear the allegedly wrongful behavior could not reasonably be expected to recur.''
\end{quote}
\noindent --- Justice Ginsburg, for the Court, 528 U.S.\ at 190.

\begin{quote}
``The relevant showing for purposes of Article III standing \ldots{} is not injury to the environment but injury to the plaintiff.''
\end{quote}
\noindent --- Justice Ginsburg, for the Court, 528 U.S.\ at 181.

\begin{quote}
``To the extent that they encourage defendants to discontinue current violations and deter them from committing future ones, they afford redress to citizen plaintiffs who are injured or threatened with injury as a consequence of ongoing unlawful conduct.''
\end{quote}
\noindent --- Justice Ginsburg, for the Court, 528 U.S.\ at 185--186.

\end{document}
